\documentclass[11pt,a4paper]{article}

%================================================
% ENCODAGE ET LANGUE
%================================================

\usepackage[utf8]{inputenc}
\usepackage[T1]{fontenc}
\usepackage[french]{babel}

%================================================
% MISE EN PAGE
%================================================

\usepackage{geometry}
\geometry{margin=1.55cm}

%================================================
% PACKAGE GTOLI
%================================================

\usepackage{../styles/gtoli}

%================================================
% TITRE DU DOCUMENT
%================================================

\title{\textbf{Titre du document}}
\author{GToli}
\date{\today}

%================================================
% DOCUMENT
%================================================

\begin{document}

\maketitle

%================================================
\section{Introduction}
%================================================

\begin{cadrebleu}{Présentation}
Ce document constitue le template de base du système pédagogique GTOLI.
\end{cadrebleu}

%================================================
\section{Exemple de rappel théorique}
%================================================

\begin{cadretheorie}
Le rappel théorique permet de mettre en évidence une notion essentielle.
\end{cadretheorie}

%================================================
\section{Exemple de méthode}
%================================================

\begin{cadremethode}
Cette section permet de présenter un protocole ou une démarche.
\end{cadremethode}

%================================================
\section{Exemple d’énoncé}
%================================================

\begin{cadreenonce}
Voici un exemple d’énoncé ou de consigne pédagogique.
\end{cadreenonce}

%================================================
\section{Exemple d’aide}
%================================================

\begin{cadreaide}
Cette section peut contenir une piste de résolution ou une aide.
\end{cadreaide}

%================================================
\section{Exemple de solution}
%================================================

\begin{cadresolution}
Voici un exemple de solution détaillée.
\end{cadresolution}

%================================================
\section{Exemple de synthèse}
%================================================

\begin{cadresynthese}
Cette section permet de résumer les éléments importants.
\end{cadresynthese}

\end{document}